% English translation of chapter2.tex lines 1245--1563.
% Source: Methods of Algebra, Volume 2, commit
% 9a5803ff77dd3257484cb177f851a73770a59dd3 (CC BY 4.0).
% stable-unit-id: o014.aljabr2.chapter2.exact-functors-injective-projective-objects
% This segment stops exactly before the source section at line 1564.

% segment-id: o014.li2.chapter2.exact-functors-injective-projective-objects.h001
\section{Exact Functors, Injective Objects, and Projective Objects}\label{sec:inj-proj}
\hypertarget{o014.aljabr2.chapter2.exact-functors-injective-projective-objects}{ }

% segment-id: o014.li2.chapter2.exact-functors-injective-projective-objects.p001
Given a functor $F: \mathcal{A} \to \mathcal{B}$, one may ask whether it
preserves $\varinjlim$ or $\varprojlim$; see \cite[\S 2.8]{Li1} or the
discussion in \S\ref{sec:limit-functor}. This section focuses on the case in
which $\mathcal{A}$ and $\mathcal{B}$ are abelian categories and $F$ is an
additive functor. Taking the special instances of $\varprojlim$ (or
$\varinjlim$) given by $\Ker$ (or $\Coker$), every morphism $f:X\to Y$ in
$\mathcal{A}$ yields a canonical morphism $F\Ker(f)\to\Ker F(f)$ and its dual
$\Coker F(f)\to F\Coker(f)$; they make the following diagrams commute.

% segment-id: o014.li2.chapter2.exact-functors-injective-projective-objects.g001
\[\begin{tikzcd}[column sep=small]
	F \Ker(f) \arrow[rd, "{F[\Ker(f) \hookrightarrow X]}"'] \arrow[rr] & & \Ker F(f) \arrow[hookrightarrow, ld] \\
	& FX &
\end{tikzcd} \quad \begin{tikzcd}[column sep=small]
	\Coker F(f) \arrow[rr] & & F \Coker(f) \\
	& FY \arrow[twoheadrightarrow, lu] \arrow[ru, "{ F[Y \twoheadrightarrow \Coker(f)] }"'] . &
\end{tikzcd}\]

% segment-id: o014.li2.chapter2.exact-functors-injective-projective-objects.l001
\begin{itemize}
% segment-id: o014.li2.chapter2.exact-functors-injective-projective-objects.i001
	\item If $F\Ker(f)\rightiso\Ker F(f)$ for every $f$, we say that
	\emph{$F$ preserves kernels};
% segment-id: o014.li2.chapter2.exact-functors-injective-projective-objects.i002
	\item if $\Coker F(f)\rightiso F\Coker(f)$ for every $f$, we say that
	\emph{$F$ preserves cokernels}.
\end{itemize}

% segment-id: o014.li2.chapter2.exact-functors-injective-projective-objects.p002
If $(X^\bullet,d^\bullet)$ is a complex in an abelian category, then
$(FX^\bullet,Fd^\bullet)$ is also a complex; the issue to be examined is
exactness.

% segment-id: o014.li2.chapter2.exact-functors-injective-projective-objects.pr001
\begin{proposition}\label{prop:left-right-exact-functor}
For an additive functor $F:\mathcal{A}\to\mathcal{B}$ between abelian
categories, the following statements are equivalent:
% segment-id: o014.li2.chapter2.exact-functors-injective-projective-objects.l002
\begin{enumerate}[(L1)]
% segment-id: o014.li2.chapter2.exact-functors-injective-projective-objects.i003
	\item $F$ preserves kernels;
% segment-id: o014.li2.chapter2.exact-functors-injective-projective-objects.i004
	\item if $0\to X'\xrightarrow{f}X\xrightarrow{g}X''$ is exact in
	$\mathcal{A}$, then
	$0\to F(X')\xrightarrow{Ff}F(X)\xrightarrow{Fg}F(X'')$ is exact in
	$\mathcal{B}$;
% segment-id: o014.li2.chapter2.exact-functors-injective-projective-objects.i005
	\item $F$ preserves all finite $\varprojlim$.
\end{enumerate}
Dually, the following statements are also equivalent:
% segment-id: o014.li2.chapter2.exact-functors-injective-projective-objects.l003
\begin{enumerate}[(R1)]
% segment-id: o014.li2.chapter2.exact-functors-injective-projective-objects.i006
	\item $F$ preserves cokernels;
% segment-id: o014.li2.chapter2.exact-functors-injective-projective-objects.i007
	\item if $X'\xrightarrow{f}X\xrightarrow{g}X''\to0$ is exact in
	$\mathcal{A}$, then
	$F(X')\xrightarrow{Ff}F(X)\xrightarrow{Fg}F(X'')\to0$ is exact in
	$\mathcal{B}$;
% segment-id: o014.li2.chapter2.exact-functors-injective-projective-objects.i008
	\item $F$ preserves all finite $\varinjlim$.
\end{enumerate}
\end{proposition}

% segment-id: o014.li2.chapter2.exact-functors-injective-projective-objects.pf001
\begin{proof}
By duality, it suffices to discuss (L1)--(L3).

(L1) $\implies$ (L2). An exact sequence of the form
$0\to\bullet\to\bullet\to\bullet$ is precisely a kernel diagram of a
morphism; see the explanation at the end of \S\ref{sec:cohomology}.

(L2) $\implies$ (L3). Every finite $\varprojlim$ can be constructed from
finite products and equalizers. We already know that additive functors preserve
biproducts; since $F$ preserves $\Ker$, it preserves all equalizers
(Remark~\ref{rem:equalizer-Ker}). Thus $F$ preserves all finite
$\varprojlim$.

(L3) $\implies$ (L1) is immediate.
\end{proof}

% segment-id: o014.li2.chapter2.exact-functors-injective-projective-objects.p003
For an additive functor $F$ as above, if every exact sequence
$(X^\bullet,d^\bullet)$ gives an exact sequence $(FX^\bullet,Fd^\bullet)$, we
say that $F$ preserves exact sequences. Likewise, one may ask whether $F$
preserves short exact sequences; these two properties turn out to be
equivalent.

% segment-id: o014.li2.chapter2.exact-functors-injective-projective-objects.pr002
\begin{proposition}\label{prop:exact-functor}
For an additive functor $F:\mathcal{A}\to\mathcal{B}$ between abelian
categories, the following statements are equivalent:
% segment-id: o014.li2.chapter2.exact-functors-injective-projective-objects.l004
\begin{enumerate}[(E1)]
% segment-id: o014.li2.chapter2.exact-functors-injective-projective-objects.i009
	\item $F$ preserves short exact sequences;
% segment-id: o014.li2.chapter2.exact-functors-injective-projective-objects.i010
	\item $F$ preserves exact sequences;
% segment-id: o014.li2.chapter2.exact-functors-injective-projective-objects.i011
	\item $F$ preserves all finite $\varprojlim$ and finite $\varinjlim$.
\end{enumerate}
\end{proposition}

% segment-id: o014.li2.chapter2.exact-functors-injective-projective-objects.pf002
\begin{proof}
(E1) $\implies$ (E2). Take an exact sequence $(X^\bullet,d^\bullet)$ in
$\mathcal{A}$ and decompose the complex into short exact sequences
% segment-id: o014.li2.chapter2.exact-functors-injective-projective-objects.g002
\[\begin{tikzcd}
	0 \arrow[r] & \Ker(d^n) \arrow[r, "{\iota^n}"] & X^n \arrow[r, "{\mathbf{d}^n}"] & \Ker(d^{n+1}) \arrow[r] & 0
\end{tikzcd}\]
where the morphism $\mathbf{d}^n$ is characterized by
$d^n=\iota^{n+1}\mathbf{d}^n$; the index $n$ may be arbitrary, except that
$X^n$ cannot be the rightmost term of the exact sequence.

By assumption, the sequence
$0\to F(\Ker(d^n))\xrightarrow{F\iota^n}F(X^n)
\xrightarrow{F\mathbf{d}^n}F(\Ker(d^{n+1}))\to0$ remains exact, and
$Fd^n=F\iota^{n+1}F\mathbf{d}^n$. The short exact sequences can therefore be
reassembled into the exact sequence $(FX^\bullet,Fd^\bullet)$ in
$\mathcal{B}$.

(E2) $\implies$ (E3). If $F$ preserves exact sequences, it satisfies
conditions (L2) and (R2) of
Proposition~\ref{prop:left-right-exact-functor}.

(E3) $\implies$ (E1). Suppose that
$0\to X'\to X\to X''\to0$ is exact in $\mathcal{A}$. Applying
Proposition~\ref{prop:left-right-exact-functor} once more, we see that
$0\to F(X')\to F(X)\to F(X'')$ and
$F(X')\to F(X)\to F(X'')\to0$ are both exact; hence
$0\to F(X')\to F(X)\to F(X'')\to0$ is exact.
\end{proof}

% segment-id: o014.li2.chapter2.exact-functors-injective-projective-objects.p004
Notice that (E3) $\iff$ (L3) $\wedge$ (R3).
% segment-id: o014.li2.chapter2.exact-functors-injective-projective-objects.d001
\begin{definition}[Exact functors]\label{def:exact-functor}
	\index{left exact, right exact, exact}
For an additive functor $F:\mathcal{A}\to\mathcal{B}$ between abelian
categories, consider conditions (L1)--(L3), (R1)--(R3) of
Proposition~\ref{prop:left-right-exact-functor}, and conditions (E1)--(E3) of
Proposition~\ref{prop:exact-functor}.
% segment-id: o014.li2.chapter2.exact-functors-injective-projective-objects.l005
\begin{itemize}
% segment-id: o014.li2.chapter2.exact-functors-injective-projective-objects.i012
	\item If any of (L1)--(L3) holds, then $F$ is called \emph{left exact}.
% segment-id: o014.li2.chapter2.exact-functors-injective-projective-objects.i013
	\item If any of (R1)--(R3) holds, then $F$ is called \emph{right exact}.
% segment-id: o014.li2.chapter2.exact-functors-injective-projective-objects.i014
	\item If any of (E1)--(E3) holds, then $F$ is called \emph{exact}; this is
	equivalent to saying that $F$ is both left exact and right exact.
\end{itemize}
\end{definition}

% segment-id: o014.li2.chapter2.exact-functors-injective-projective-objects.p005
For example, every equivalence between abelian categories is of course an
exact functor. Another extreme example is the zero functor: it sends every
object to a zero object and every morphism to a zero morphism; this functor is
also exact.

% segment-id: o014.li2.chapter2.exact-functors-injective-projective-objects.p006
By Lemma~\ref{prop:mono-epi-ker-coker}, a left exact functor preserves
monomorphisms, while a right exact functor preserves epimorphisms.

% segment-id: o014.li2.chapter2.exact-functors-injective-projective-objects.rem001
\begin{remark}\label{rem:exactness-mono-epi}
In view of (E1), if $F$ is known to be left exact (or right exact), then $F$ is
exact if and only if $F$ preserves epimorphisms (or monomorphisms).
\end{remark}

% segment-id: o014.li2.chapter2.exact-functors-injective-projective-objects.p007
A standard way to check left or right exactness is by means of adjoint
functors.

% segment-id: o014.li2.chapter2.exact-functors-injective-projective-objects.th001
\begin{theorem}\label{prop:exactness-adjoint}
Consider a pair of additive functors between abelian categories $\mathcal{A}$
and $\mathcal{B}$:
% segment-id: o014.li2.chapter2.exact-functors-injective-projective-objects.g003
\[\begin{tikzcd}
	F: \mathcal{A} \arrow[r, yshift=0.3em] & \mathcal{B} :G \arrow[l, yshift=-0.3em] .
\end{tikzcd}\]
Suppose that $F$ and $G$ can be equipped as an adjoint pair $(F,G,\varphi)$.
Then $F$ is right exact and $G$ is left exact; indeed, $F$ preserves
$\varinjlim$ and $G$ preserves $\varprojlim$.
\end{theorem}

% segment-id: o014.li2.chapter2.exact-functors-injective-projective-objects.pf003
\begin{proof}
Apply \cite[Theorem 2.8.12]{Li1} and conditions (L3), (R3) of
Proposition~\ref{prop:left-right-exact-functor}.
\end{proof}

% segment-id: o014.li2.chapter2.exact-functors-injective-projective-objects.ex001
\begin{example}[Left/right exactness of limits]\label{eg:limit-exactness}
Let $\mathcal{A}$ be an abelian category, let $I$ be any category, and suppose
that every functor $\alpha:I\to\mathcal{A}$ has a $\varinjlim$ (or a
$\varprojlim$). Endow $\mathcal{A}^I$ with the abelian-category structure of
Proposition~\ref{prop:functor-cat-Abel}. We shall show that the functor
$\varinjlim:\mathcal{A}^I\to\mathcal{A}$ (or $\varprojlim$) is right exact (or
left exact).

It suffices to discuss $\varinjlim$. By
Theorem~\ref{prop:exactness-adjoint}, one strategy is to show that
$\varinjlim$ has a right adjoint. Define the diagonal functor
$\Delta:\mathcal{A}\to\mathcal{A}^I$, which sends $L\in\Obj(\mathcal{A})$ to
the constant functor $\Obj(I)\ni i\mapsto L$. We have an adjoint pair
% segment-id: o014.li2.chapter2.exact-functors-injective-projective-objects.g004
\[\begin{tikzcd}
	\varinjlim : \mathcal{A}^I \arrow[shift left, r] & \mathcal{A}: \Delta \arrow[l, shift left] .
\end{tikzcd}\]
This follows immediately from the universal property of $\varinjlim$: in
$\mathcal{A}^I$, giving a morphism from a functor
$\alpha:I\to\mathcal{A}$ to $\Delta(L)$ is equivalent to giving a compatible
family of morphisms $f_i:\alpha(i)\to L$; in other words, a cone with base
$\alpha$ and vertex $L$. See the related review in
\S\ref{sec:limit-functor}.
\end{example}

% segment-id: o014.li2.chapter2.exact-functors-injective-projective-objects.ex002
\begin{example}\label{eg:module-adjoint-pairs}
Let $f:R\to S$ be a ring homomorphism. Consider the following functors between
categories of left modules:
% segment-id: o014.li2.chapter2.exact-functors-injective-projective-objects.g005
\[\begin{tikzcd}[column sep=large]
	R\dcate{Mod} \arrow[bend left=50, r, "{S \dotimes{R} (\cdot)}"] \arrow[bend right=50, r, "{\Hom_R({}_R S, \cdot)}"'] & S\dcate{Mod} \arrow[l, "{ {}_{R \to S} \mathcal{F} }"]
\end{tikzcd}\]
Here the forgetful functor ${}_{R\to S}\mathcal{F}$ is simply the operation of
regarding an $S$-module as an $R$-module via $f$. The functors
$S\dotimes{R}(\cdot)$ and $\Hom_R({}_R S,\cdot)$ are discussed in
\cite[\S 6.6]{Li1}; here ${}_R S$ means $S$ regarded as a left $R$-module.
According to \cite[Corollary 6.6.8]{Li1},
% segment-id: o014.li2.chapter2.exact-functors-injective-projective-objects.q001
\[\left(S\dotimes{R}-,\;{}_{R\to S}\mathcal{F}\right)
\quad\text{and}\quad
\left({}_{R\to S}\mathcal{F},\;\Hom_R({}_R S,-)\right)\]
are adjoint pairs. Theorem~\ref{prop:exactness-adjoint} then implies
% segment-id: o014.li2.chapter2.exact-functors-injective-projective-objects.q002
\[S\dotimes{R}(\cdot)\ \text{is right exact},\quad
\Hom({}_R S,\cdot)\ \text{is left exact},\quad
{}_{R\to S}\mathcal{F}\ \text{is exact}.\]
This exactness can also be checked directly by algebraic means. For example,
for the forgetful functor ${}_{R\to S}\mathcal{F}$, whether a sequence of
module homomorphisms
$\cdots\to M^n\xrightarrow{d^n}M^{n+1}\to\cdots$ is a complex (that is,
$d^{n+1}d^n=0$), or is exact (that is,
$\Image(d^n)=\Ker(d^{n+1})$), depends only on the additive-group structure of
$M^n$, not on multiplication in $R$ or $S$. In other words,
${}_{\Z\to S}\mathcal{F}:S\dcate{Mod}\to\cate{Ab}$ is already an exact
functor. Alternatively, one can check directly that
${}_{R\to S}\mathcal{F}$ preserves all limits; this follows directly from
the construction in \cite[Theorem 6.2.2]{Li1}.
\end{example}
% segment-id: o014.li2.chapter2.exact-functors-injective-projective-objects.p008
An exact functor automatically preserves the $\Image\simeq\Coim$ of a
morphism (Definition~\ref{def:Im-Coim},
Proposition~\ref{prop:Im-Coim-additive}); it also preserves cohomology.

% segment-id: o014.li2.chapter2.exact-functors-injective-projective-objects.pr003
\begin{proposition}\label{prop:exact-preserves-cohomologies}
Let $F:\mathcal{A}\to\mathcal{B}$ be an exact functor between abelian
categories. For every complex $(X^\bullet,d^\bullet)$ in $\mathcal{A}$, there
is a canonical isomorphism in $\mathcal{B}$
% segment-id: o014.li2.chapter2.exact-functors-injective-projective-objects.q003
\[F\Hm^n(X^\bullet,d^\bullet)\rightiso
\Hm^n(FX^\bullet,Fd^\bullet),\qquad n\in\Z.\]
\end{proposition}

% segment-id: o014.li2.chapter2.exact-functors-injective-projective-objects.pf004
\begin{proof}
It suffices to consider a three-term complex
$X'\xrightarrow{f}X\xrightarrow{g}X''$. Returning to definition
\eqref{eqn:homology-3},
% segment-id: o014.li2.chapter2.exact-functors-injective-projective-objects.q004
\begin{multline*}
	F\left(\Hm\left[X'\xrightarrow{f}X\xrightarrow{g}X''\right]\right)
	=F\Coker\left[\Image(f)\to\Ker(g)\right]\\
	\simeq\Coker\left[F\Image(f)\to F\Ker(g)\right]\\
	\simeq\Coker\left[\Image(Ff)\to\Ker(Fg)\right]
	=\Hm\left[FX'\xrightarrow{Ff}FX\xrightarrow{Fg}FX''\right],
\end{multline*}
and all the isomorphisms that occur are canonical.
\end{proof}

% segment-id: o014.li2.chapter2.exact-functors-injective-projective-objects.p009
Faithful exact functors have particularly useful properties.

% segment-id: o014.li2.chapter2.exact-functors-injective-projective-objects.pr004
\begin{proposition}\label{prop:faithful-exact}\index{faithfully exact functor}
For an additive functor $F:\mathcal{A}\to\mathcal{B}$ between abelian
categories, the following statements are equivalent:
% segment-id: o014.li2.chapter2.exact-functors-injective-projective-objects.l006
\begin{enumerate}[(i)]
% segment-id: o014.li2.chapter2.exact-functors-injective-projective-objects.i015
	\item $F$ is exact and faithful;
% segment-id: o014.li2.chapter2.exact-functors-injective-projective-objects.i016
	\item $F$ is exact, and $FX=0\iff X=0$ for every
	$X\in\Obj(\mathcal{A})$;
% segment-id: o014.li2.chapter2.exact-functors-injective-projective-objects.i017
	\item $X'\to X\to X''$ is exact in $\mathcal{A}$ if and only if
	$FX'\to FX\to FX''$ is exact in $\mathcal{B}$.
\end{enumerate}
\end{proposition}

% segment-id: o014.li2.chapter2.exact-functors-injective-projective-objects.pf005
\begin{proof}
(i) $\implies$ (ii): If $FX=0$, then
$F(\identity_X)=\identity_{FX}=0$, whence $\identity_X=0$ and $X=0$.

(ii) $\implies$ (iii): Apply
Proposition~\ref{prop:exact-preserves-cohomologies} to the cohomology.

(iii) $\implies$ (i): This condition immediately implies that $F$ is exact.
Suppose that $u:X\to Y$ satisfies $Fu=0$. Since
% segment-id: o014.li2.chapter2.exact-functors-injective-projective-objects.q005
\[X\xrightarrow{\identity}X\xrightarrow{u}Y\xrightarrow{\identity}Y,\]
the image of this sequence under $F$ is exact; hence the sequence itself is
exact and $u=0$.
\end{proof}

% segment-id: o014.li2.chapter2.exact-functors-injective-projective-objects.p010
Localization (see \S\ref{sec:cat-localization}) provides an important example
of an exact functor.

% segment-id: o014.li2.chapter2.exact-functors-injective-projective-objects.pr005
\begin{proposition}[Exactness of localization]\label{prop:Abel-cat-localization}
Let $\mathcal{A}$ be an abelian category and let
$S\subset\Mor(\mathcal{A})$ be a multiplicative family
(Definition~\ref{def:multiplicative-family}). Then the category
$\mathcal{A}[S^{-1}]$ supplied by Theorem~\ref{prop:Gabriel-Zisman} has a
canonical abelian-category structure, and the localization functor
$Q:\mathcal{A}\to\mathcal{A}[S^{-1}]$ is exact.
\end{proposition}

% segment-id: o014.li2.chapter2.exact-functors-injective-projective-objects.pf006
\begin{proof}
Theorem~\ref{prop:localization-additivity} gives
$\mathcal{A}[S^{-1}]$ a canonical additive-category structure under which $Q$
is an additive functor. We shall show that every morphism $f$ in
$\mathcal{A}[S^{-1}]$ has a kernel and cokernel and is strict. By the
construction of $\mathcal{A}[S^{-1}]$, after composing $f$ with an isomorphism
from $S$, we may assume that $f$ is the image under $Q$ of a morphism in
$\mathcal{A}$; this does not change the property to be proved.
Lemma~\ref{prop:localization-prod} says that $Q$ preserves all finite
$\varinjlim$ and $\varprojlim$. In particular, $Q$ sends $\Ker$, $\Coker$,
$\Image$, and $\Coim$ in $\mathcal{A}$ to the corresponding constructions in
$\mathcal{A}[S^{-1}]$, and therefore also preserves diagram
\eqref{eqn:strict-morphism}. This shows that $\mathcal{A}[S^{-1}]$ is an
abelian category; condition (E3) of Proposition~\ref{prop:exact-functor} shows
that $Q$ is exact.
\end{proof}

% segment-id: o014.li2.chapter2.exact-functors-injective-projective-objects.p011
If, in addition, $\mathcal{A}$ is a $\Bbbk$-linear abelian category, where
$\Bbbk$ is a commutative ring, then $Q$ is a functor between $\Bbbk$-linear
abelian categories. This too is part of
Theorem~\ref{prop:localization-additivity}.

% segment-id: o014.li2.chapter2.exact-functors-injective-projective-objects.p012
Another extremely important example is the $\Hom$ functor. If $T$ is an
object of an abelian category $\mathcal{A}$, then $\Hom(T,\cdot)$ defines a
functor $\mathcal{A}\to\cate{Ab}$, while $\Hom(\cdot,T)$ defines a functor
$\mathcal{A}^{\opp}\to\cate{Ab}$; on morphisms, they send $f:X\to Y$ to
$f_*$ and $f^*$ on $\Hom$, respectively. Both are plainly additive functors.

% segment-id: o014.li2.chapter2.exact-functors-injective-projective-objects.p013
If $\mathcal{A}$ is a $\Bbbk$-linear abelian category, the $\Hom$ functors may
take values in $\Bbbk\dcate{Mod}$ and become $\Bbbk$-linear. This extension
has little effect on the discussion below, so we do not treat it separately.

% segment-id: o014.li2.chapter2.exact-functors-injective-projective-objects.pr006
\begin{proposition}[Left exactness of $\Hom$]\label{prop:Hom-left-exact}
Let $\mathcal{A}$ be an abelian category and let $T$ be an object of
$\mathcal{A}$. Then $\Hom(T,\cdot):\mathcal{A}\to\cate{Ab}$ and
$\Hom(\cdot,T):\mathcal{A}^{\opp}\to\cate{Ab}$ are both left exact functors.
\end{proposition}

% segment-id: o014.li2.chapter2.exact-functors-injective-projective-objects.pf007
\begin{proof}
By duality (replacing $\mathcal{A}$ by $\mathcal{A}^{\opp}$), it suffices to
treat $\Hom(T,\cdot)$. The point is to show that $\Hom(T,\cdot)$ preserves
$\Ker$. Since $\Ker$ in $\cate{Ab}$ is simply the kernel from group theory,
this follows immediately from the universal property of $\Ker$.
\end{proof}

% segment-id: o014.li2.chapter2.exact-functors-injective-projective-objects.d002
\begin{definition}\label{def:injective-projective-obj}
	\index{injective object}
	\index{projective object}
Let $X$ be an object of an abelian category $\mathcal{A}$. If
$\Hom(X,\cdot):\mathcal{A}\to\cate{Ab}$ is an exact functor, then $X$ is
called a \emph{projective object}; if
$\Hom(\cdot,X):\mathcal{A}^{\opp}\to\cate{Ab}$ is an exact functor, then $X$
is called an \emph{injective object}.
\end{definition}

% segment-id: o014.li2.chapter2.exact-functors-injective-projective-objects.p014
By Remark~\ref{rem:exactness-mono-epi} and
Proposition~\ref{prop:Hom-left-exact}, to decide whether an object $X$ is
projective (or injective), it suffices to check whether the functor
$\Hom(X,\cdot)$ (or $\Hom(\cdot,X)$) preserves epimorphisms. Consequently:
% segment-id: o014.li2.chapter2.exact-functors-injective-projective-objects.l007
\begin{itemize}
% segment-id: o014.li2.chapter2.exact-functors-injective-projective-objects.i018
	\item An object $P$ is projective if and only if, for every exact sequence
	in $\mathcal{A}$ of the form $Y\xrightarrow{g}X\to0$ and every morphism
	$P\to X$, there exists a $P\to Y$ making the following diagram commute.
% segment-id: o014.li2.chapter2.exact-functors-injective-projective-objects.g006
	\[\begin{tikzcd}
		& P \arrow[d] \arrow[ld, "\exists"'] & \\
	Y \arrow[r, "g"'] & X \arrow[r] & 0 \\
	\end{tikzcd}\]
	(Equivalently, $g_*:\Hom(P,Y)\to\Hom(P,X)$ is surjective.)
% segment-id: o014.li2.chapter2.exact-functors-injective-projective-objects.i019
	\item An object $I$ is injective if and only if, for every exact sequence
	in $\mathcal{A}$ of the form $0\to X\xrightarrow{f}Y$ and every morphism
	$X\to I$, there exists a $Y\to I$ making the following diagram commute.
% segment-id: o014.li2.chapter2.exact-functors-injective-projective-objects.g007
	\[\begin{tikzcd}
		0 \arrow[r] & X \arrow[r, "f"] \arrow[d] & Y \arrow[ld, "\exists"] \\
		& I &
	\end{tikzcd}\]
	(Equivalently, $f^*:\Hom(Y,I)\to\Hom(X,I)$ is surjective.)
\end{itemize}

% segment-id: o014.li2.chapter2.exact-functors-injective-projective-objects.lm001
\begin{lemma}\label{prop:inj-proj-split}
Consider a short exact sequence
$0\to X'\xrightarrow{f}X\xrightarrow{g}X''\to0$ in an abelian category. If
$X'$ is injective or $X''$ is projective, then this short exact sequence
splits.
\end{lemma}

% segment-id: o014.li2.chapter2.exact-functors-injective-projective-objects.pf008
\begin{proof}
By duality, we may assume without loss of generality that $X'$ is injective.
Apply the characterization above to the commutative diagram
% segment-id: o014.li2.chapter2.exact-functors-injective-projective-objects.g008
\[\begin{tikzcd}
	0 \arrow[r] & X' \arrow[r, "f"] \arrow[d, "{\identity_{X'}}"'] & X \arrow[ld, "\exists r"] \\
	& X' &
\end{tikzcd}\]
The relation $rf=\identity_{X'}$ then satisfies the condition of
Proposition~\ref{prop:split-ses}.
\end{proof}

% segment-id: o014.li2.chapter2.exact-functors-injective-projective-objects.lm002
\begin{lemma}\label{prop:prod-injective-projective}
Consider a family of objects $(X_i)_{i\in I}$ in an abelian category
$\mathcal{A}$. If the coproduct $\coprod_{i\in I}X_i$ (or the product
$\prod_{i\in I}X_i$) exists in $\mathcal{A}$, then the coproduct is projective
(or the product is injective) if and only if every $X_i$ is so.
\end{lemma}

% segment-id: o014.li2.chapter2.exact-functors-injective-projective-objects.pf009
\begin{proof}
By duality, it suffices to consider the coproduct $\coprod_{i\in I}X_i$. The
universal property gives an isomorphism of functors
% segment-id: o014.li2.chapter2.exact-functors-injective-projective-objects.q006
\[
\Hom_{\mathcal{A}}\left(\coprod_{i\in I}X_i,-\right)
\rightiso\prod_{i\in I}\Hom_{\mathcal{A}}\left(X_i,-\right):
\mathcal{A}\to\cate{Ab}.
\]
It remains only to use the following elementary observation: for a family of
	maps $(f_i:A_i\to B_i)_{i\in I}$, where $A_i,B_i$ are sets, the induced map
$(f_i)_{i\in I}:\prod_{i\in I}A_i\to\prod_{i\in I}B_i$ is surjective if and
only if every $f_i$ is surjective.
\end{proof}

% segment-id: o014.li2.chapter2.exact-functors-injective-projective-objects.p015
For example, let $R$ be a ring. Every free $R$-module is a projective object
of $R\dcate{Mod}$. Indeed, it suffices to show that $R$ itself is projective;
but $\Hom(R,\cdot):R\dcate{Mod}\to\cate{Ab}$ is isomorphic to the forgetful
functor $\mathcal{F}:R\dcate{Mod}\to\cate{Ab}$ via the isomorphism that sends
a homomorphism $\varphi:R\to X$ to $\varphi(1)\in X$. Thus
$\Hom(R,\cdot)$ is exact.

% segment-id: o014.li2.chapter2.exact-functors-injective-projective-objects.p016
Free modules are also generators for $R\dcate{Mod}$; see
Definition~\ref{def:generators} and Example~\ref{eg:generators}. An injective
object that is also a cogenerator (or a projective object that is also a
generator) is very useful. Here is a simple characterization.

% segment-id: o014.li2.chapter2.exact-functors-injective-projective-objects.pr007
\begin{proposition}[Injective cogenerators and projective generators]\label{prop:injective-cogenerator}
In an abelian category $\mathcal{A}$, an injective (or projective) object $X$
is a cogenerator (or generator) if and only if
$\Hom(T,X)\neq0$ (or $\Hom(X,T)\neq0$) for every
$T\in\Obj(\mathcal{A})$ with $T\neq0$.

Equivalently, $\Hom(\cdot,X)$ (or $\Hom(X,\cdot)$) is a faithful exact
functor.
\end{proposition}

% segment-id: o014.li2.chapter2.exact-functors-injective-projective-objects.pf010
\begin{proof}
We treat only the case in which $X$ is injective. First suppose that $X$ is a
cogenerator. Apply the definition of cogenerator to
% segment-id: o014.li2.chapter2.exact-functors-injective-projective-objects.g009
\[\begin{tikzcd}
	T \arrow[r, shift left, "{\identity_T}"] \arrow[r, shift right, "0"'] & T
\end{tikzcd}\]
There is then a $\delta\in\Hom(T,X)$ such that
$\delta=\delta\circ\identity_T\neq\delta\circ0=0$.

Conversely, we wish to show that if $h:S\to T$ is nonzero, then there is a
$\delta\in\Hom(T,X)$ such that $\delta h\neq0$. Observe that there is a
$\delta'\in\Hom(\Image(h),X)$ with $\delta'\neq0$. Since
$S\twoheadrightarrow\Image(h)$ is an epimorphism, the composite
$S\twoheadrightarrow\Image(h)\xrightarrow{\delta'}X$ is nonzero. Now use the
injectivity of $X$ (see the discussion following
Definition~\ref{def:injective-projective-obj}) to extend $\delta'$ to
$\delta:T\to X$.

The assertion about faithful exact functors is an immediate application of
Proposition~\ref{prop:faithful-exact}.
\end{proof}

% segment-id: o014.li2.chapter2.exact-functors-injective-projective-objects.p017
To do homological algebra in an abelian category, one often needs enough
injective or projective objects.

% segment-id: o014.li2.chapter2.exact-functors-injective-projective-objects.d003
\begin{definition}\label{def:enough-injectives}
	\index{enough injectives}
	\index{enough projectives}
Let $\mathcal{A}$ be an abelian category.
% segment-id: o014.li2.chapter2.exact-functors-injective-projective-objects.l008
\begin{itemize}
% segment-id: o014.li2.chapter2.exact-functors-injective-projective-objects.i020
	\item If, for every $X\in\Obj(\mathcal{A})$, there are an injective object
	$I$ and a monomorphism $X\hookrightarrow I$, then $\mathcal{A}$ is said to
	\emph{have enough injectives}.
% segment-id: o014.li2.chapter2.exact-functors-injective-projective-objects.i021
	\item If, for every $X\in\Obj(\mathcal{A})$, there are a projective object
	$P$ and an epimorphism $P\twoheadrightarrow X$, then $\mathcal{A}$ is said
	to \emph{have enough projectives}.
\end{itemize}
\end{definition}

% segment-id: o014.li2.chapter2.exact-functors-injective-projective-objects.p018
These two notions are dual to each other. A standard way to construct
injective or projective objects is to use adjoints of exact functors; here are
the details.

% segment-id: o014.li2.chapter2.exact-functors-injective-projective-objects.pr008
\begin{proposition}\label{prop:adjoint-injective-projective}
Consider a pair of functors between abelian categories
% segment-id: o014.li2.chapter2.exact-functors-injective-projective-objects.g010
\[\begin{tikzcd}
	\mathcal{A} \arrow[r, shift left, "F"] & \mathcal{B} \arrow[l, shift left, "G"]
\end{tikzcd}\]
and suppose that $F$ is exact.
% segment-id: o014.li2.chapter2.exact-functors-injective-projective-objects.l009
\begin{itemize}
% segment-id: o014.li2.chapter2.exact-functors-injective-projective-objects.i022
	\item If $G$ is a left adjoint of $F$, then $G$ sends projective objects in
	$\mathcal{B}$ to projective objects in $\mathcal{A}$;
% segment-id: o014.li2.chapter2.exact-functors-injective-projective-objects.i023
	\item if $G$ is a right adjoint of $F$, then $G$ sends injective objects in
	$\mathcal{B}$ to injective objects in $\mathcal{A}$.
\end{itemize}
\end{proposition}

% segment-id: o014.li2.chapter2.exact-functors-injective-projective-objects.pf011
\begin{proof}
By duality, it suffices to consider the case in which $G$ is a left adjoint.
In this case $G$ is necessarily additive
(Corollary~\ref{prop:automatic-additivity}). Let $P$ be a projective object of
$\mathcal{B}$. For every exact sequence $(X^\bullet,d^\bullet)$ in
$\mathcal{A}$, there is an isomorphism of complexes in $\cate{Ab}$
% segment-id: o014.li2.chapter2.exact-functors-injective-projective-objects.q007
\[
\Hom_{\mathcal{A}}(GP,X^\bullet)\rightiso
\Hom_{\mathcal{B}}(P,FX^\bullet).
\]
Since $F$ is exact, $(FX^\bullet,Fd^\bullet)$ is an exact sequence, and hence
the right-hand side is exact in $\cate{Ab}$. Therefore
$\Hom_{\mathcal{A}}(GP,\cdot):\mathcal{A}\to\cate{Ab}$ is an exact functor.
\end{proof}

% segment-id: o014.li2.chapter2.exact-functors-injective-projective-objects.p019
For example, consider a ring $R$ and the forgetful functor
$\mathcal{F}:R\dcate{Mod}\to\cate{Ab}$. By
Example~\ref{eg:module-adjoint-pairs}, $\mathcal{F}$ is exact and has the right
adjoint $G:=\Hom_{\cate{Ab}}(R,\cdot)$.
Proposition~\ref{prop:adjoint-injective-projective} shows that $G$ sends
injective objects in $\cate{Ab}$ to injective objects in $R\dcate{Mod}$.
Injective objects in $\cate{Ab}$ are easy to characterize: they are precisely
the divisible $\Z$-modules. This is the standard way to construct injective
$R$-modules in module theory and to show that $R\dcate{Mod}$ has enough
injectives; see \cite[Theorem 6.9.14]{Li1}.

% segment-id: o014.li2.chapter2.exact-functors-injective-projective-objects.p020
On the other hand, $R\dcate{Mod}$ also has enough projectives: for every
$R$-module $X$, choose a set of generators $A\subset X$; then
$R^{\oplus A}\twoheadrightarrow X$.

% segment-id: o014.li2.chapter2.exact-functors-injective-projective-objects.ex003
\begin{example}\label{eg:A2-injective-projective}
The following combined exercise involves an abstract abelian category and will
be used in \S\sourcecrossref{sec:derived-primer}{3.12} to study long exact
sequences of derived functors. Let $\mathcal{A}$ be an abelian category.
Consider the category $\mathbf{2}$ (whose diagram is $0\to1$; see the book's
introduction on categories). The functor category
$\mathcal{A}^{\mathbf{2}}$ is again abelian
(Proposition~\ref{prop:functor-cat-Abel}); it is the arrow category: its
objects are morphisms $X_0\to X_1$ in $\mathcal{A}$, and its morphisms are
commutative squares in $\mathcal{A}$
% segment-id: o014.li2.chapter2.exact-functors-injective-projective-objects.g011
\[\begin{tikzcd}[row sep=small, column sep=small]
	X_0 \arrow[r] \arrow[d] & Y_0 \arrow[d] \\
	X_1 \arrow[r] & Y_1
\end{tikzcd}\]

% segment-id: o014.li2.chapter2.exact-functors-injective-projective-objects.p021
For $i\in\{0,1\}$, the evaluation functor
$\mathrm{ev}_i:\mathcal{A}^{\mathbf{2}}\to\mathcal{A}$ sends an object
$X_0\to X_1$ to $X_i$. In addition, define additive functors from
$\mathcal{A}$ to $\mathcal{A}^{\mathbf{2}}$ by
\footnote{Editorial correction (O014-C027): the source writes
$H:X\to[X\xrightarrow{\identity_X}X]$, whereas the context defines the object
map of the functor $H$, in parallel with $L_0$ and $L_1$. The assignment arrow
$\mapsto$ is used here.}
% segment-id: o014.li2.chapter2.exact-functors-injective-projective-objects.q008
\[L_0:X\mapsto[X\to0],\quad L_1:X\mapsto[0\to X],\quad
H:X\mapsto[X\xrightarrow{\identity_X}X],\]
for $X\in\Obj(\mathcal{A})$; the definitions on morphisms are clear. We shall
prove the following assertions.
% segment-id: o014.li2.chapter2.exact-functors-injective-projective-objects.l010
\begin{compactenum}[(i)]
% segment-id: o014.li2.chapter2.exact-functors-injective-projective-objects.i024
	\item The functors $\mathrm{ev}_0$ and $\mathrm{ev}_1$ are both exact and
	both send injective (or projective) objects of
	$\mathcal{A}^{\mathbf{2}}$ to injective (or projective) objects of
	$\mathcal{A}$. Moreover, $H$, $L_0$, and $L_1$ are all exact.
% segment-id: o014.li2.chapter2.exact-functors-injective-projective-objects.i025
	\item The functors $L_0,H$ (or $H,L_1$) send injective (or projective)
	objects of $\mathcal{A}$ to injective (or projective) objects of
	$\mathcal{A}^{\mathbf{2}}$.
% segment-id: o014.li2.chapter2.exact-functors-injective-projective-objects.i026
	\item If $\mathcal{A}$ has enough injectives (or projectives), then so does
	$\mathcal{A}^{\mathbf{2}}$.
\end{compactenum}

% segment-id: o014.li2.chapter2.exact-functors-injective-projective-objects.p022
Since $\varinjlim$ and $\varprojlim$ in $\mathcal{A}^{\mathbf{2}}$ are
constructed objectwise, the functors $\mathrm{ev}_0,\mathrm{ev}_1$ and
$H,L_0,L_1$ are indeed exact. The remainder of (i) and (ii) follows from the
adjoint relations satisfied by $\mathrm{ev}_i$; the diagrams are as follows,
and verification is a simple exercise:
% segment-id: o014.li2.chapter2.exact-functors-injective-projective-objects.g012
\[\begin{tikzcd}[column sep=huge]
	\mathcal{A}^{\mathbf{2}} \arrow[r, "{\mathrm{ev}_0}" description] & \mathcal{A} \arrow[l, bend right, "{H: \text{left adjoint}}"'] \arrow[l, bend left, "L_0: \text{right adjoint}"] 
\end{tikzcd} \quad \begin{tikzcd}[column sep=huge]
	\mathcal{A}^{\mathbf{2}} \arrow[r, "{\mathrm{ev}_1}" description] & \mathcal{A} \arrow[l, bend right, "{L_1: \text{left adjoint}}"'] \arrow[l, bend left, "{H: \text{right adjoint}}"] .
\end{tikzcd}\]

% segment-id: o014.li2.chapter2.exact-functors-injective-projective-objects.p023
For (iii), first consider injective objects. Take an object
$[X_0\xrightarrow{f}X_1]$ of $\mathcal{A}^{\mathbf{2}}$, and choose
monomorphisms $\epsilon_i:X_i\hookrightarrow I_i$, where $I_i$ is injective,
$i\in\{0,1\}$. These give two morphisms in
$\mathcal{A}^{\mathbf{2}}$, which may be written as the commutative diagrams
% segment-id: o014.li2.chapter2.exact-functors-injective-projective-objects.g013
\[\begin{tikzcd}[column sep=large, /tikz/execute at end picture={
		\node[rectangle, draw, dashed, fit=(NE) (SE), inner xsep=0.5em, inner ysep=0em] {};
	}]
	X_0 \arrow[d, "f"'] \arrow[hookrightarrow, r, "\epsilon_0"] & |[alias=NE]| I_0 \arrow[d] \\
	X_1 \arrow[r] & |[alias=SE]| 0
\end{tikzcd} \quad \text{and} \quad \begin{tikzcd}[column sep=large, /tikz/execute at end picture={
	\node[rectangle, draw, dashed, fit=(NE) (SE), inner xsep=0.5em, inner ysep=0em] {};
	}]
	X_0 \arrow[d, "f"'] \arrow[r, "\epsilon_1 f"] & |[alias=NE]| I_1 \arrow[d, "\identity"] \\
	X_1 \arrow[hookrightarrow, r, "\epsilon_1"'] & |[alias=SE]| I_1
\end{tikzcd}\]
The two framed columns are $L_0(I_0)$ and $H(I_1)$, respectively. By (ii),
both are injective objects in $\mathcal{A}^{\mathbf{2}}$, and so is their
direct sum (Lemma~\ref{prop:prod-injective-projective}). Hence the commutative
diagram
% segment-id: o014.li2.chapter2.exact-functors-injective-projective-objects.g014
\[\begin{tikzcd}[column sep=large]
	X_0 \arrow[d, "f"'] \arrow[hookrightarrow, r, "{(\epsilon_0, \epsilon_1 f)}"] & I_0 \oplus I_1 \arrow[d, "\text{projection}"] \\
	X_1 \arrow[hookrightarrow, r, "{\epsilon_1}"'] & I_1
\end{tikzcd}\]
embeds $[X_0\xrightarrow{f}X_1]$ into an injective object. For projective
objects, use the functors $H$ and $L_1$.

% segment-id: o014.li2.chapter2.exact-functors-injective-projective-objects.p024
\end{example}

% segment-id: o014.li2.chapter2.exact-functors-injective-projective-objects.p025
The discussion in Example~\ref{eg:A2-injective-projective} extends from
$\mathcal{A}^{\mathbf{2}}$ to a general functor category
$\mathcal{A}^{\mathcal{C}}$; the argument presents no essentially new
difficulty. The exercises in this chapter give the details.
